\section{Image Processing And Plant Phenotyping}

 \paragraph{}
 Traditional plant phenotyping methods are destructive, prone to human error, time-consuming, and challenging to perform from start to end of a plant's ontogenesis \cite{arnalbarbedoDigitalImageProcessing2013, zhangHighthroughputPhenotypingPlant2023}. Recent advances in image sensing, processing, and data management have produced high-throughput image-based plant phenotyping systems, which are amenable to automation and can analyse large samples of data precisely in a very short period. Reducing the need for potentially erroneous, time-consuming, expensive human labour in plant phenotyping \cite{daschoudhuryLeveragingImageAnalysis2019}. 
 
 \paragraph{}
 These Modern high-throughput image-based plant phenotyping systems are examples of cyberinfrastructure (CI) systems \cite{abebeImageBasedHighThroughputPhenotyping2023}, which are research environments that support data acquisition, storage, management and computing and processing services distributed over some network \cite{nabwireReviewApplicationArtificial2021}. For image-based plant phenotyping, this means that image acquisition and processing techniques are vital.
 
 \subsection{Image Acquistion Techniques In Plant Phenotyping Systems}
 
 \paragraph{}
  
Image processing techniques in plant phenotyping cannot be applied until suitable images are acquired. Image sensor advancements have provided the foundations for detecting various information from plants using electromagnetic radiation. Plant tissues interact with this radiation by absorbing, reflecting, emitting, transmitting, and fluorescing light. Different ranges of wavelengths correspond to different biochemical or physiological aspects of plant tissue, which are not visible to the naked eye \cite{liReviewImagingTechniques2014}. Thus, the choice of imaging technique for data acquisition affects the types of phenotypical parameters that can be measured. 

\subparagraph{}
Current imaging techniques are visible, fluorescence, thermal, tomographic, and hyperspectral imaging \cite{SensorsMeasuringPlant, liReviewImagingTechniques2014}, where the most commonly used techniques are visible, and hyperspectral imaging due to their affordablity, robustness and flexibility in measuring a wide range of phenotypic parameters.

\subparagraph{}
The most common imaging technique is visible-imaging. Typically afforadable, and high-resolution RGB cameras are used in visible imaging. These devices are sensitive to the visible spectral range and have been used to measure various phenotypic parameters like shoot biomass, yield and leaf morphology, and salinity tolerance \cite{golzarianAccurateInferenceShoot2011, hoyos-villegasGroundBasedDigitalImaging2014}.
 
 \subparagraph{}
 Hyperspectral imaging is another widely used in plant phenotyping. In this technique, sensors capture multiple spectral ranges, enabling detailed measurements for various phenotypic parameters. Notable applications of hyperspectral imaging include estimating nutrient content and detecting plant responses to abiotic or biotic stressors \cite{saricApplicationsHyperspectralImaging2022}.
 
 \subparagraph{}
The other imaging techniques are not as widely used as visible or hyperspectral imaging; however, they are also effective in measuring plant phenotypic parameters. For instance, thermal imaging in plant water stress relations \cite{liReviewImagingTechniques2014, ishimweApplicationsThermalImaging2014}., fluorescence imaging in used plant pathogenic symptom monitoring \cite{lichtenthalerRoleChlorophyllFluorescence1988, swarbrickMetabolicConsequencesSusceptibility2006}, and tomography in estimating plant organ water content \cite{windtMRILongdistanceWater2006}.

\subsection{Image Processing Techniques In Plant Phenotyping Systems}
There is a lack of generalized or standardised image processing techniques in plant phenotyping; hence, bespoke image processing pipelines are designed for specific plant species or genotypes. However,  many image-processing pipelines in plant phenotyping perform segmentation, machine learning or deep-learning systems for automated analysis and processing of phenotypic data.

 
 % Below is for draft